% 第21章 复核清单
\section{复核清单：已执行结果与剩余待办}\label{sec:checklist}

这一章把本综述做过的核验、跑过的实验、以及尚未完成的部分如实列出，供读者判断本文每条结论的可靠程度。

\subsection{检索与来源池}

\begin{itemize}
  \item \textbf{多源检索已执行}：OpenAlex 标题检索（覆盖资产定价、最优执行、强化学习、NLP、生成式、图网络、评测方法论等锚点文献）、arXiv 关键词检索（q-fin 与 cs 分类，覆盖六个方法族与五个交易环节）、Crossref 补核验。
  \item \textbf{来源池规模}：去重后共 528 条候选来源（期刊/Crossref 178 条、arXiv 350 条）；目录 \texttt{outputs/pool.json}。检索脚本见 \texttt{scripts/retrieve\_sources.py}。
  \item \textbf{引用键核验}：正文引用的来源逐条经 OpenAlex/Crossref/arXiv 核对元数据，脚本见 \texttt{scripts/verify\_refs.py}，结果见 \texttt{outputs/refs\_verified.json}。
\end{itemize}

\subsection{参考实现实跑记录（P0--P2）}

\texttt{code/run.py} 在合成时点面板（$T=1200$，$N=60$，$K=8$，快照校验和 \texttt{432bdae19c9533ed}）上端到端运行，候选池 61 个（8 个单特征 + 50 个随机线性组合 + 3 个岭参数网格），试验账本口径见下。关键输出如下：

\begin{table}[htbp]
\centering
\caption{参考实现实跑结果（合成数据，见下方性质声明）}
\label{tab:runresult}
\begin{tabular}{@{}>{\raggedright\arraybackslash}p{5.6cm}>{\raggedright\arraybackslash}p{8.8cm}@{}}
\toprule
项目 & 实跑结果 \\
\midrule
不可成交单元占比 & 0.1921（含涨跌停 / 停牌 / 非成分） \\
试验账本 $N$ & 61（单特征 8 + 随机组合 50 + 岭参数网格 3） \\
样本内最优候选 & \texttt{rand11}（\textbf{纯噪声}线性组合），样本内夏普 1.196 \\
其滚动前向 OOS 夏普 & 0.477（单次切分为 0.868，两者不一致本身即是警示） \\
OOS 夏普高于它的纯噪声候选数 & 15 \\
信号类指标（该候选） & IC 均值 0.0074、IR 0.052、IC$>$0 占比 0.501（与纯噪声一致） \\
成本三档（该候选，全样本） & 零成本 夏普 0.833 / 固定费率 0.518 / 含冲击 0.057 \\
缩水夏普 DSR（$N=61$） & 标准误口径 DSR $=0.053$（$\text{SR}_0=0.096$）；截面离散口径 DSR $=0.027$ \\
回测过拟合概率 PBO & 0.567；切分敏感性 $\{8{:}0.400,\ 12{:}0.567,\ 16{:}0.516,\ 20{:}0.584\}$，区间 $[0.400, 0.584]$ \\
\bottomrule
\end{tabular}
\end{table}

这组数字本身就是本文论点的演示：\textbf{按样本内挑选，选出来的是一条纯噪声组合；它的 IC 与零无异；扣掉冲击成本后夏普几乎归零；而 DSR 与 PBO 都明确指出，这个结果经不起选择偏差与过拟合的检验。}流程的价值不在于跑出收益，而在于它能在结果看起来不错时就先给出否定信号。

\noindent\textbf{性质声明（重要）}：本演示使用\textbf{合成面板}，其信噪比、涨跌停比例、停牌率均为作者设定。因此它\textbf{只能验证评测流程本身是否按预期工作，不能作为“真实市场中普遍存在过拟合”的证据}——在一个低信噪比的设定下，“按样本内选择会选中噪声”几乎是被设计出来的结果。要把它上升为经验结论，必须在公开真实数据上做同样的最小实证；这一步已在下一小节补上。

还有一处口径差异需注明：本演示的成本用\textbf{单一 $5$bp 费率与常数 ADV $10^9$ 元}（合成面板不含成交额字段，无法逐股取 ADV），而下一小节的真实数据实证用\textbf{佣金＋过户费＋分段印花税与逐股 ADV}。两处的“成本三档”因此\textbf{不可直接比较}，只能各自内部比较。

\subsection{真实数据最小实证（P0--P2）}\label{sec:real-empirical}

合成面板只能验证流程本身，不能替代经验证据。为补上这一步，本文令同一套 P0--P2 流程直接跑在\textbf{本地通达信 vipdoc 日线}（\texttt{.day}）上。复现命令 \texttt{python code/run.py {-}{-}real}，数据加载见 \texttt{code/real\_data.py}。\footnote{本节的口径与举例在修订过程中更正过四处：冲击成本的 ADV 由常数改为逐股取值、封板占比统一为有效单元格分母、复权核验由抽样扩到全量 400 只、单特征异象举例按多重检验结果更正。更正前后的对照与影响见 \texttt{docs/change\_summary\_round2.tex}，正文只保留更正后的口径。}

\textbf{数据与口径}：样本为沪深（含北交所 15 只）A 股，窗口 2019-01-02 至 2026-10-08（$T=1881$ 个交易日）。universe 要求窗口内交易日覆盖率不低于 90\%（按上市后活动窗口计）、日成交额中位数不低于 2000 万元，按代码等距抽样至 $N=400$。抽样框依赖本地库清单，故\textbf{随快照一并冻结}于 \texttt{outputs/real\_universe.json}（400 只代码与逐股 ADV），他人可用同一清单复算。

价格底座为通达信原始日线，再乘后复权因子，先于一切特征与收益计算，消除分红与送转造成的跳空。这里有一处须如实说明的外部依赖：通达信本地权息文件 \texttt{gbbq} 在本机三个安装中均为\textbf{加密存储}，无法解析，故复权因子取 AKShare 并缓存至 \texttt{outputs/hfq\_factors.json}（第三方仅作必要兜底）。该拼装链已对\textbf{全部 400 只逐股交叉核验}：通达信原始价乘因子与 AKShare 独立后复权价相比，相对偏差中位数 $2.3\times10^{-4}$、95 分位 $9.1\times10^{-4}$、最大 $3.1\times10^{-3}$（17 只超过 $10^{-3}$，无一只超过 $10^{-2}$）。核验脚本 \texttt{scripts/verify\_adjust.py}，结果缓存于 \texttt{outputs/adjust\_check.json}。

交易制度掩码按\textbf{封板价精确判定}：对 $\{5\%,10\%,20\%,30\%\}$ 四档涨跌幅分别计算 $\text{round}(\text{前收}\times(1\pm\text{lim}),0.01)$，当收盘价与最高价（涨停）或最低价（跌停）同时等于该价时判为封板、不可调仓。该判定逐日直接识别，ST 的 $\pm5\%$ 与主板/创业科创/北交所的 $10/20/30\%$ 自动区分，不依赖名称表、也不引入前视。一字板（开$=$高$=$低$=$收且封板）与停牌一并记为不可成交；新股上市后前 20 个交易日剔除。以下占比一律\textbf{以有效单元格为分母}（有效单元格指已上市且特征可算的单元，占全部单元的 80.1\%）：封板合计 1.96\%，其中涨停 1.49\%、跌停 0.48\%；一字板 0.20\%。

\textbf{成本口径}：冲击成本按\textbf{逐个股 ADV} 计算——ADV 取该股在窗口内日成交额的中位数（样本中位为 $7.46\times10^7$ 元），参与率 $=$ 成交额 $/$ ADV，冲击率 $=0.02\sqrt{\text{参与率}}$，参与率设 10\% 上限。费率按 A 股实际构成计入：佣金 2.5bp 与过户费 0.1bp 双边，另加\textbf{卖出印花税分段}（2023-08-28 之前 10bp、之后 5bp）。作为对照：若对全部个股取同一常数 ADV（如 $10^9$ 元，恰为样本中位数的 13.4 倍），参与率会被系统性低估——这是把成本算轻最常见的一种做法，本文不使用。

特征仍是 $K=8$ 的价量因子（1/5/20/60 日动量、20 日波动、量比、距 20 日均线偏离、Amihud 非流动性），逐期截面 $z$ 标准化，只用 $t$ 期及以前信息；标签为后复权收盘到收盘的下一日收益。快照校验和 \texttt{781cadd68dfbd170}。

\begin{table}[htbp]
\centering
\caption{参考实现实跑结果（本地通达信真实日线，含后复权、交易制度掩码与逐股成本口径）}
\label{tab:runreal}
\begin{tabular}{@{}>{\raggedright\arraybackslash}p{5.6cm}>{\raggedright\arraybackslash}p{8.8cm}@{}}
\toprule
项目 & 实跑结果 \\
\midrule
数据来源 & 通达信 vipdoc 日线（价格）+ AKShare 后复权因子（第三方兜底，400 只全量核验） \\
窗口与规模 & 2019-01-02 至 2026-10-08；$T=1881$，$N=400$（含北交所 15 只、窗口内上市 153 只） \\
不可成交单元占比 & 0.2134；其中有效样本内封板 1.96\%，其余为上市前与历史不足 \\
试验账本 $N$ & 61（单特征 8 + 随机组合 50 + 岭参数网格 3） \\
样本内最优候选 & \texttt{rand10}（随机线性组合），样本内夏普 3.649 \\
其滚动前向 OOS 夏普 & 2.671（单次切分 3.595） \\
OOS 夏普高于它的候选数 & 1 \\
信号类指标（该候选） & IC 均值 0.0129、IR 0.088、IC$>$0 占比 0.548（$n=1820$） \\
换手率 & 日均合计 $\sum|\Delta w|=0.721$，年化约 182 倍 \\
成本三档（年化 / 夏普） & 零成本 0.273 / 3.057；固定费率 0.154 / 1.720；含冲击 $-0.134$ / $-1.492$ \\
成本占毛收益 & 零成本 0；固定费率 0.44；含冲击 1.49（即成本已超过毛收益本身） \\
缩水夏普 DSR（$N=61$） & 标准误口径 0.9991；截面离散口径 0.0022（取舍见下文） \\
回测过拟合概率 PBO & 0.000；切分敏感性区间 $[0.000, 0.000]$ \\
\bottomrule
\end{tabular}
\end{table}

\textbf{这组数字推翻了合成演示的一条外推。}在合成面板上，“按样本内挑选会选中噪声”几乎是设定好的结果；到真实数据上，同样按样本内夏普挑选，选中的随机组合样本内 3.649、滚动前向 OOS 2.671，DSR 标准误口径 0.9991、PBO 0.000，都指向“不是选择偏差的产物”。\textbf{合成面板得到的“选到噪声”不能外推为经验规律}——这正是最小实证该做的事：用真实数据检验框架自己。

按第 \ref{sec:repro_protocol} 章自己援引的标准，单特征结论必须过多重检验校正，故此处补齐 $t$ 值（$t=\text{IC-IR}\times\sqrt{n}$，$n=1820$）与 Benjamini--Hochberg 校正，结果见表 \ref{tab:lslo} 后四列。该式假设逐期 IC 序列独立，故另做序列相关检验：各特征 IC 的 $\text{AR}(1)$ 多在 $0.05$ 以内，Newey--West 调整后仅非流动性由 $3.92$ 降到 $3.45$（仍过门槛）、量比几乎不动，可见上述判定对序列相关稳健（脚本 \texttt{code/audit\_significance.py}）。

\textbf{DSR 两口径的取舍}：两个口径的标准化方式相同，差别只在 $sr_0$ 中 $V(\widehat{SR})$ 的取法——标准误口径用单个夏普估计的解析方差，截面离散口径用 61 个候选样本外夏普的截面方差。本文\textbf{以标准误口径为准}，理由是：\textbf{截面离散度把两类东西混在了一起}——各候选真实夏普的差异，与每个候选夏普的估计噪声。本池里真实强弱差异很大（量比与波动率的真实表现相差悬殊），实测两者的合成量比纯估计方差大 20 倍，多出来的部分主要是质量差异而非估计误差，因此它不能充当 $V(\widehat{SR})$；代入后会把 $sr_0$ 抬到 $0.2556$（标准误口径为 $0.0722$）。

另做一项稳健性检查：候选样本外收益的平均两两相关约为零（$-0.012$，最大 $0.977$），说明 $N=61$ 作为试验数是合适的、无需按相关结构折减；即便改用按相关结构折算的有效试验数（$193$）重算，$sr_0$ 也只是由 $0.0652$ 升到 $0.0764$，DSR 由 $0.9999$ 变为 $0.9995$——\textbf{结论对这一选择不敏感}。仍须如实留下两点边界：原文 DSR 的 $V(\widehat{SR})$ 取自试验集夏普的方差，解析方差是试验集不可用时的替代，两者并不等价；且 $E[\max_N]$ 项假设试验相互独立，候选相关会使其偏大，故标准误口径在该方向上略偏保守。

需要说清的是，这个“随机组合”在合成与真实两种数据上的含义并不相同。合成面板里只有 3 个特征带信号，随机组合大多由噪声主导；真实数据里 8 个价量特征本身都含信号，随机组合只是它们的随机线性混合，因此“选到噪声”这一合成结论本就没有外推基础。

复权口径本身改变了什么，本文另做了一项对照：保持特征与权重不变，只把标签由后复权换成不复权重算，8 个特征的截面 IC 变化均不超过 $0.0011$（脚本 \texttt{code/diag\_adjust\_effect.py}）。也就是说，在日频价量特征上，除权跳空并不是截面相关的显著来源——这一点与“必须复权否则信号被污染”的直觉预期相反，故如实记下。

\begin{table}[htbp]
\centering
\caption{单特征对照：多空（不可实施）与仅多头（可实施）的 OOS 夏普，以及截面 IC 的显著性}
\label{tab:lslo}
\begin{tabular}{@{}>{\raggedright\arraybackslash}p{2.1cm}>{\raggedright\arraybackslash}p{1.6cm}>{\raggedright\arraybackslash}p{1.7cm}>{\raggedright\arraybackslash}p{1.3cm}>{\raggedright\arraybackslash}p{1.3cm}>{\raggedright\arraybackslash}p{1.7cm}>{\raggedright\arraybackslash}p{4.6cm}@{}}
\toprule
特征 & 多空 OOS 夏普 & 仅多头 OOS 夏普 & IC-IR & $t$ & 通过校正 & 读法 \\
\midrule
量比 & 3.056 & 0.838 & $-0.212$ & $-9.03$ & BH、$|t|>3$ & 多空最强，方向压倒性依赖做空 \\
1 日反转 & 1.211 & 0.435 & $0.159$ & $6.77$ & BH、$|t|>3$ & 多头化后大幅回落 \\
非流动性 & 1.516 & 1.271 & $0.092$ & $3.91$ & BH、$|t|>3$ & 多头端最强，但混入市场 beta \\
20 日动量 & 1.409 & 0.694 & $-0.056$ & $-2.38$ & 仅 BH & 多头化后回落 \\
60 日动量 & 1.200 & 0.737 & $-0.053$ & $-2.28$ & 仅 BH & 同上 \\
距均线偏离 & 1.952 & 0.767 & $-0.044$ & $-1.89$ & 未通过 & 多空强、多头弱，IC 不显著 \\
波动率（低波） & 0.862 & 0.884 & $-0.040$ & $-1.71$ & 未通过 & 多头端贡献不低于多空 \\
5 日动量 & 2.142 & 0.739 & $-0.031$ & $-1.32$ & 未通过 & 多空强，但 IC 不显著 \\
\bottomrule
\end{tabular}
\end{table}

读这张表要注意两件事。第一，\textbf{“多空夏普高”与“IC 显著”并不一致}：距均线偏离与 5 日动量的多空 OOS 夏普分别有 1.952 与 2.142，但它们的截面 IC 在 BH 校正下并不显著；反过来说，真正稳健的只有 3 个特征（量比、1 日反转、非流动性，均过 $|t|>3$ 门槛），另有 2 个仅过 BH。第二，按同一标准，\textbf{5 日动量不应与量比并列}——它 $|t|=1.32$，既不过 BH 也不过 $|t|>3$；而真正过硬门槛的 1 日反转与非流动性在原文举例中反而没被提到。

表注：末列“通过校正”中，BH 指 Benjamini--Hochberg 5\% 校正，$|t|>3$ 指 Harvey--Liu--Zhu 门槛。两者严格程度并不同阶——$|t|>3$ 相当于 $p<0.003$ 的单次检验，而 BH 在多检验中段更宽松，因此上表按“两项都过 / 仅过 BH / 未通过”三级标注，判定顺序是先看 $|t|>3$ 再看 BH。

\textbf{统计上稳健，不等于可交易；但“不可交易”也不是一个笼统的结论。}门槛有三道：做空的可实施性、换手带来的成本、以及市场暴露。前两道与\textbf{换手}直接相关，而换手在候选之间的跨度极大，因此必须逐候选看，不能用一个候选代表全部——详见下一小节。

\subsection{可实施性核查：成本、换手与市场暴露}

前两小节说明“信息层确有稳定信号”。这一小节要问的是：这些信号里\textbf{哪些能扛住摩擦}。答案高度依赖候选——换手在候选之间的跨度是 26 倍，用一个候选代表全部会得出错误结论。把 8 个特征的两种构建方式各算一遍（\texttt{code/audit\_cost\_by\_candidate.py}），得到表 \ref{tab:costbycand}。

\begin{table}[htbp]
\centering
\caption{候选级换手与成本敏感性（全样本；符号按原特征方向，未取样本内更优方向）}
\label{tab:costbycand}
\begin{tabular}{@{}>{\raggedright\arraybackslash}p{3.4cm}>{\raggedright\arraybackslash}p{1.9cm}>{\raggedright\arraybackslash}p{2.1cm}>{\raggedright\arraybackslash}p{2.1cm}>{\raggedright\arraybackslash}p{4.6cm}@{}}
\toprule
特征（口径） & 年化换手 & 零成本夏普 & 含冲击夏普 & 读法 \\
\midrule
非流动性（多空） & 14 & 1.221 & \textbf{1.014} & 成本后仍成立 \\
非流动性（仅多头） & 18 & 1.026 & \textbf{0.892} & 成本后仍成立 \\
波动率（多空） & 28 & $-1.123$ & $-1.518$ & 换手低，但信号方向为负 \\
60 日动量（多空） & 46 & $-1.394$ & $-1.984$ & 中低换手 \\
20 日动量（多空） & 77 & $-1.630$ & $-2.800$ & 成本耗损约 1.2 \\
距均线偏离（多空） & 92 & $-1.998$ & $-3.496$ & 成本耗损约 1.5 \\
量比（多空） & 154 & $-3.422$ & $-7.190$ & 高换手 \\
5 日动量（多空） & 152 & $-2.150$ & $-5.166$ & 高换手 \\
1 日反转（多空） & 322 & $-1.294$ & $-10.235$ & 换手最高，成本耗损 8.9 \\
\texttt{rand10}（多空） & 182 & 3.057 & $-1.492$ & 论文表 17 所用候选 \\
\bottomrule
\end{tabular}
\end{table}

表 \ref{tab:costbycand} 里最要紧的一行是\textbf{非流动性}：年化换手只有 14--18 倍，含冲击后多空夏普仍为 $1.014$、仅多头仍为 $0.892$。16 种构建方式中只有这 2 种在含冲击后为正。换手分布为最小 14 倍、中位 86 倍、最大 358 倍。因此正确的表述不是“可预测但不可交易”，而是——\textbf{成本后能否存活取决于换手，低换手候选可以活下来，高换手候选活不下来}。表 17 所用的 \texttt{rand10} 恰落在最贵的一类里（年化换手 182 倍，含冲击档 $-1.492$、成本占毛收益 1.49），它演示的是这个条件的一端，而不是全部。

关于非流动性这一行，必须配一句限定，而且这条线索已被单独跟进（见 \texttt{docs/illiquidity\_note.tex}）。这类因子的收益很可能反映的是\textbf{流动性补偿}而非可提取的 alpha——Amihud 类信号靠持有别人不愿持有的股票获利，而成本模型只能捕捉正常时段的成交冲击，捕捉不到“需要卖的时候恰好卖不掉”这件事。跟进分析给出了量化支持：按当期 ADV 分位\textbf{时点安全地}逐级筛除较不流动的个股后，含冲击夏普单调衰减（全部个股 $1.02$、保留较流动的 70\% 为 $0.69$、50\% 为 $0.47$、30\% 仅 $0.08$），即\textbf{这份收益长在最不流动的那一端}。另有两点支撑与两点保留：十层收益近乎完全单调（层序相关 $0.975$，不是尾部噪声）、与其余七个因子的相关最大不过 $0.27$（不是已知异象的替身）、对成本假设不敏感（ADV 打到四分之一后仍有 $0.91$）；但逐年看 2019--2020 两年含冲击为负（$-1.47$、$-0.30$），稳定性不如全样本夏普暗示的那样好，且流动性压力期依赖的均值差 $t$ 值仅 $1.47$、不显著。

这一读法还要排除一个竞争解释：它会不会只是小市值因子？按单独方案做的规模中性化（流通市值取逐日股本，时点安全）给出了“无法分离”的中间答案——中性化后毛夏普由 $1.221$ 降到 $1.001$（规模解释不了全部），但含冲击夏普被压到 $0.611$，跌破预设的 $0.7$ 阈值，主因是换手从 14 倍升到 19 倍；且层内价差随规模单调衰减（最小规模层 $15.1$ 个基点对最大规模层 $2.0$ 个基点），效应确实集中在小盘端。\textbf{结论是：它过了成本关，但没过多可提取性关；且“不可提取”的理由经中性化后更清楚——收益集中在低流动端，中性化本身还会抬高换手。}

接下来是最后一道门：市场暴露。用同期等权全市场组合作基准（年化 12.1\%、夏普 0.523、滚动前向 OOS 夏普 0.597），对两个代表候选各做一元回归分解——\texttt{rand10} 与分析中存活下来的非流动性，结果见表 \ref{tab:betadecomp}。

\begin{table}[htbp]
\centering
\caption{beta / alpha 分解（基准为同期等权全市场组合；alpha 为年化）}
\label{tab:betadecomp}
\begin{tabular}{@{}>{\raggedright\arraybackslash}p{2.4cm}>{\raggedright\arraybackslash}p{2.4cm}>{\raggedright\arraybackslash}p{1.4cm}>{\raggedright\arraybackslash}p{1.9cm}>{\raggedright\arraybackslash}p{1.9cm}>{\raggedright\arraybackslash}p{4.0cm}@{}}
\toprule
候选（口径） & beta & alpha & OOS 夏普 & 剥 beta 后夏普 & 读法 \\
\midrule
\texttt{rand10}（多空） & 0.025 & 27.0\% & 2.671 & 3.030 & 近似市场中性 \\
\texttt{rand10}（仅多头） & 0.991 & 5.5\% & 0.780 & 1.082 & 几乎全是市场暴露 \\
非流动性（多空） & 0.027 & 11.4\% & 1.516 & 1.190 & 近似市场中性 \\
非流动性（仅多头） & 1.005 & 16.6\% & 1.271 & 1.061 & 市场暴露高，但 alpha 厚 \\
基准本身 & 1.000 & —— & 0.597 & —— & 等权全市场 \\
\bottomrule
\end{tabular}
\end{table}

这张表把两个候选的差别摆清楚了。\texttt{rand10} 的仅多头口径 beta 为 $0.991$，基本是“加了一层选股的市场组合”，剥离市场暴露后年化 alpha 只剩 $5.5\%$；而非流动性虽然 beta 同样接近 1（$1.005$），剥离后的年化 alpha 有 $16.6\%$——比 \texttt{rand10} 厚三倍。两者叠加换手差异（182 倍对 18 倍），结论就统一起来了：\textbf{可交易性不是“信号强弱”的函数，而是“换手与 alpha 之比”的函数}。表 18 里“仅多头 OOS 夏普 1.271”这类数字，其 beta 部分确实很大，但它对应的 alpha 也并非可以忽略；而 \texttt{rand10} 的 $0.780$ 则主要是 beta。

还有一处必须交代：样本内最优的 \texttt{rand10} 虽是“随机线性组合”，但它与\textbf{量比的逐期截面相关为 $-0.862$}——它基本是量比因子的反向近似（另与 1 日反转相关 $-0.545$、与非流动性 $-0.270$）。所以“随机组合在真实数据上胜出”并不意味着凭空出现了一个新信号，而是它把最强的单因子反着复制了一遍再稀释。这一点是对合成演示那一节的必要补充。

这组核查落在第 \ref{sec:foundations} 章的缝隙判据上：信息（IC 显著、DSR 通过）只是必要条件。从“可预测”到“可交易”，中间横着做空可实施性、换手成本与市场暴露三道门；真实数据把这句论断从口号变成了可量化的落差，并且给出了落差量的分布——而不是一个统一的否定。按本节与跟进分析的结果，三道门的通过情况可以分三档：\texttt{rand10} 这类高换手候选在第一道门就被成本拦下；非流动性过了成本关却卡在可提取性上（收益长在最不流动的一端，且有两个年度为负）；而真正三道全过的候选，在这 8 个因子里\textbf{一个也没有}。

\noindent\textbf{边界（重要，且方向上偏乐观）}：universe 只保留窗口末仍在市、且成交额达标的股票，存在\textbf{幸存者偏差}；作为基准的等权全市场组合是日频、零成本、且同样建在存活样本上，alpha 与 beta 均继承了这一口径的偏差；复权因子来自第三方而非本地权息文件，全量核验下最大相对偏差 $3.1\times10^{-3}$，若该因子在某处滞后于真实除权日仍会残留局部跳空；涨跌停判定按四档固定幅度尝试，未处理上市首日、重大资产重组复牌等无涨跌幅限制或特殊安排；ADV 用窗口内成交额中位数而非滚动值，会平滑掉流动性骤降的时段；未按停牌频率进一步筛除标的。这些都会影响具体数值，方向上倾向于\textbf{高估可交易性}，故上文的可实施性对照应读作一个上界性质的结果，而非收益预测。


\subsection{核验状态约定}

证据表中的每条来源标注核验状态。本文实际使用两级：\textbf{已核验}（DOI/arXiv/官方链接与元数据互证）、\textbf{部分核验}（元数据确认，样本区间或参数细节待全文复核）。若仅能定位来源、无法确证元数据，则不入表。本文对定性叙述广泛引用，对具体数值一律谨慎——凡未在正文给出精确数值的结论，均因原文数值未经全文逐项核验。

\subsection{剩余待办与已知缺口}

\begin{enumerate}
  \item \textbf{参数级复核}：证据表中标注“部分核验/待复核”的条目，其样本区间、阈值、成本模型等细节仍需逐篇读全文核对。
  \item \textbf{复现状态的实际验证}：第 \ref{sec:repro_audit} 章的复现状态核验表依据公开仓库的可见信息（可用性、许可、基准披露、活跃度）判定；对“可复现”等级的工作，尚未全部做端到端实跑复算。
  \item \textbf{私有数据结论不可验证}：另类数据与商业 LLM 类结论，因数据/模型不可得，本文只能标注“不可独立验证”。
  \item \textbf{非英语文献}：检索以英文为主，中文及其他语种的实证工作可能被低估。
  \item \textbf{真实数据实证的边界待收敛}：第 \ref{sec:real-empirical} 节已补后复权、交易制度掩码、逐股成本口径、全量复权核验、候选级换手对照与 beta 分解，但复权因子依赖第三方（本地权息文件加密不可解析），universe 带幸存者偏差，ADV 用窗口内中位数而非滚动值，且成本后的存活判定只在“多空／仅多头 × 8 个单特征”这 16 种构建方式上做过——更细的权重构造（如分层、行业中性）未逐一检验。
  \item \textbf{非流动性跟进的未决项}：规模中性化已完成并给出“无法分离”的结论（见 \texttt{docs/illiquidity\_note.tex}），但股本数据只有单一来源（本地 \texttt{base.dbf} 布局不可解析、无法交叉校验）。已对该来源做过五项形态自检（全部通过）与两次外部核对，残余风险仅限规模中性化一节、且方向使结论偏乐观；详见 \texttt{docs/share\_source\_note.tex}。压力期检验的证据也偏弱（$t=1.47$）。
\end{enumerate}

\noindent 本文的态度是：\textbf{一份综述的可信度，不在于它覆盖了多少，而在于它诚实标注了多少没做到。}
